\paragraph{稳定层二态统计极限与规范体积修正的一阶账本}
在稳定层均匀经验谱的两点弱极限、统一层矩函数以及 fold-gauge 常数恢复链已经建立之后，还可把二态一致渐近对均匀层平均量与规范体积修正的支配压缩为下列四条直接后果。它们分别给出缩放退化度的量化两点律、全部缩放矩与幂和首项常数、均匀平均对数退化度的二态展开，以及 \(g_m-\kappa_m+1\) 的精确首项。

\begin{theorem}[均匀稳定类型下缩放退化度的量化两点极限律]\label{thm:xi-time-part70a-uniform-scaled-degeneracy-quantitative}
令
\[
u_m(w):=\frac{1}{|X_m|}\qquad (w\in X_m),
\]
并取 \(W_m\sim u_m\)。定义缩放退化度随机变量
\[
Y_m:=\left(\frac{\varphi}{2}\right)^m d_m^{\mathrm{bin}}(W_m).
\]
则对任意有界 Lipschitz 函数 \(f:\RR\to\RR\)，都有
\[
\mathbb E[f(Y_m)]
=
\frac{F_{m+1}}{F_{m+2}}\,f(1)
+
\frac{F_m}{F_{m+2}}\,f(\varphi^{-1})
+O\!\left(\left(\frac{\varphi}{2}\right)^m\right),
\]
其中隐含常数仅依赖于 \(f\) 的 Lipschitz 常数。特别地，\(Y_m\) 依分布收敛到两点随机变量
\[
Y_\infty=
\begin{cases}
1,&\text{概率 } \varphi^{-1},\\[3pt]
\varphi^{-1},&\text{概率 } \varphi^{-2}.
\end{cases}
\]
\end{theorem}
\begin{proof}
由定理 \ref{thm:fold-bin-two-state-asymptotic}，存在常数 \(C>0\)，使得对充分大的 \(m\) 与全部 \(w\in X_m\) 都一致有
\[
\left|
\left(\frac{\varphi}{2}\right)^m d_m^{\mathrm{bin}}(w)-\varphi^{-w_m}
\right|
\le C\left(\frac{\varphi}{2}\right)^m.
\]
因此
\[
\bigl|Y_m-\varphi^{-W_m(m)}\bigr|
\le C\left(\frac{\varphi}{2}\right)^m
\]
几乎处处成立。若记 \(L_f\) 为 \(f\) 的 Lipschitz 常数，则
\[
\bigl|f(Y_m)-f(\varphi^{-W_m(m)})\bigr|
\le
L_f C\left(\frac{\varphi}{2}\right)^m.
\]
对均匀分布取期望，得到
\[
\mathbb E[f(Y_m)]
=
\mathbb E[f(\varphi^{-W_m(m)})]
+O\!\left(\left(\frac{\varphi}{2}\right)^m\right).
\]
另一方面，
\[
\mathbb E[f(\varphi^{-W_m(m)})]
=
\frac{F_{m+1}}{F_{m+2}}\,f(1)
+
\frac{F_m}{F_{m+2}}\,f(\varphi^{-1}),
\]
因为
\[
\mathbb P(W_m(m)=0)=\frac{F_{m+1}}{F_{m+2}},
\qquad
\mathbb P(W_m(m)=1)=\frac{F_m}{F_{m+2}}.
\]
这就得到所示定量公式。再用
\[
\frac{F_{m+1}}{F_{m+2}}\to\varphi^{-1},
\qquad
\frac{F_m}{F_{m+2}}\to\varphi^{-2},
\]
即可得到所述两点极限分布。证毕。
\end{proof}

\begin{corollary}[统一缩放矩极限与 \texorpdfstring{$S_q^{\mathrm{bin}}(m)$}{Sqbin(m)} 的首项常数]\label{cor:xi-time-part70a-scaled-moments-leading-constant}
对任意实数 \(q\ge 0\)，都有
\[
\lim_{m\to\infty}
\frac{1}{|X_m|}
\sum_{w\in X_m}
\left(
\left(\frac{\varphi}{2}\right)^m d_m^{\mathrm{bin}}(w)
\right)^q
=
\frac{1}{\varphi}+\varphi^{-q-2}.
\]
等价地，若记
\[
S_q^{\mathrm{bin}}(m):=\sum_{w\in X_m}\bigl(d_m^{\mathrm{bin}}(w)\bigr)^q,
\]
则有精确首项渐近
\[
\boxed{\;
S_q^{\mathrm{bin}}(m)
\sim
\frac{\varphi+\varphi^{-q}}{\sqrt5}
\left(\frac{2^q}{\varphi^{q-1}}\right)^m.\;}
\]
\end{corollary}
\begin{proof}
对固定 \(q\ge 0\)，由上一条定理作用于函数 \(f_q(y):=y^q\) 即可。事实上，由定理 \ref{thm:xi-time-part70a-uniform-scaled-degeneracy-quantitative} 的证明，对充分大的 \(m\) 与几乎处处的样本点都有
\[
Y_m\in\left[\frac{\varphi^{-1}}{2},2\right],
\]
故 \(f_q\) 在该固定紧区间上 Lipschitz。于是
\[
\frac{1}{|X_m|}
\sum_{w\in X_m}
\left(
\left(\frac{\varphi}{2}\right)^m d_m^{\mathrm{bin}}(w)
\right)^q
\to
\frac{1}{\varphi}+\varphi^{-q-2}.
\]
将上式改写为
\[
S_q^{\mathrm{bin}}(m)
=
|X_m|
\left(\frac{2}{\varphi}\right)^{mq}
\left(\frac{1}{\varphi}+\varphi^{-q-2}+o(1)\right),
\]
再用 Binet 公式
\[
|X_m|=F_{m+2}\sim \frac{\varphi^{m+2}}{\sqrt5},
\]
便得到
\[
S_q^{\mathrm{bin}}(m)
\sim
\frac{\varphi^{m+2}}{\sqrt5}
\left(\frac{2}{\varphi}\right)^{mq}
\left(\varphi^{-1}+\varphi^{-q-2}\right)
=
\frac{\varphi+\varphi^{-q}}{\sqrt5}
\left(\frac{2^q}{\varphi^{q-1}}\right)^m.
\]
证毕。
\end{proof}

\begin{theorem}[均匀稳定层平均对数退化度的二态展开]\label{thm:xi-time-part70a-uniform-average-logdeg-two-state}
定义均匀稳定层平均对数退化度
\[
\overline{\ell}_m
:=
\frac{1}{|X_m|}\sum_{w\in X_m}\log d_m^{\mathrm{bin}}(w).
\]
则有
\[
\boxed{\;
\overline{\ell}_m
=
m\log\!\Bigl(\frac{2}{\varphi}\Bigr)
-\frac{F_m}{F_{m+2}}\log\varphi
+O\!\left(\left(\frac{\varphi}{2}\right)^m\right).\;}
\]
因而
\[
\boxed{\;
\overline{\ell}_m
=
m\log\!\Bigl(\frac{2}{\varphi}\Bigr)
-\varphi^{-2}\log\varphi
+O\!\left(\left(\frac{\varphi}{2}\right)^m\right).\;}
\]
\end{theorem}
\begin{proof}
由定理 \ref{thm:fold-bin-two-state-asymptotic} 的一致对数展开，
\[
\log d_m^{\mathrm{bin}}(w)
=
m\log\!\Bigl(\frac{2}{\varphi}\Bigr)
-w_m\log\varphi
+O\!\left(\left(\frac{\varphi}{2}\right)^m\right)
\qquad (w\in X_m),
\]
且误差在 \(w\) 上一致。对全部 \(w\in X_m\) 求平均，得到
\[
\overline{\ell}_m
=
m\log\!\Bigl(\frac{2}{\varphi}\Bigr)
-\frac{1}{|X_m|}\sum_{w\in X_m}w_m\,\log\varphi
+O\!\left(\left(\frac{\varphi}{2}\right)^m\right).
\]
而
\[
\sum_{w\in X_m}w_m=|A_1|=F_m,
\qquad
|X_m|=F_{m+2},
\]
故
\[
\frac{1}{|X_m|}\sum_{w\in X_m}w_m=\frac{F_m}{F_{m+2}}.
\]
代回即得第一式。再用 Binet 公式给出的比值极限
\[
\frac{F_m}{F_{m+2}}=\varphi^{-2}+O(\varphi^{-2m}),
\]
便得到第二式。证毕。
\end{proof}
\leanverified{paper\_xi\_time\_part70a\_uniform\_average\_logdeg\_two\_state}

\leanverified{paper\_xi\_time\_part70a\_gauge\_fiber\_gap\_first\_order}
\begin{theorem}[规范体积与纤维信息差的一阶首项]\label{thm:xi-time-part70a-gauge-fiber-gap-first-order}
定义
\[
\mu_m(w):=\frac{d_m^{\mathrm{bin}}(w)}{2^m},
\qquad
\kappa_m:=\sum_{w\in X_m}\mu_m(w)\log d_m^{\mathrm{bin}}(w),
\qquad
g_m:=2^{-m}\log|G_m^{\mathrm{bin}}|,
\]
并沿用推论 \ref{cor:fold-bin-recover-2pi} 的规范常数
\[
C_m
:=
\exp\!\left(
\frac{2^{m+1}}{|X_m|}\bigl(g_m-\kappa_m+1\bigr)-\overline{\ell}_m
\right).
\]
则有
\[
\boxed{\;
\frac{2^{m+1}}{|X_m|}\bigl(g_m-\kappa_m+1\bigr)
=
m\log\!\Bigl(\frac{2}{\varphi}\Bigr)
-\frac{F_m}{F_{m+2}}\log\varphi
+\log(2\pi)
+O\!\left(\left(\frac{\varphi}{2}\right)^m\right).\;}
\]
因此
\[
\boxed{\;
g_m-\kappa_m+1
=
\frac{\varphi^2}{2\sqrt5}
\Bigl(
m\log\!\Bigl(\frac{2}{\varphi}\Bigr)
-\varphi^{-2}\log\varphi
+\log(2\pi)
\Bigr)
\left(\frac{\varphi}{2}\right)^m
+O\!\left(\left(\frac{\varphi}{2}\right)^{2m}\right).\;}
\]
\end{theorem}
\begin{proof}
由 \(C_m\) 的定义可得精确恒等式
\[
\frac{2^{m+1}}{|X_m|}\bigl(g_m-\kappa_m+1\bigr)=\log C_m+\overline{\ell}_m.
\]
另一方面，推论 \ref{cor:fold-bin-recover-2pi} 给出
\[
\log C_m=\log(2\pi)+O\!\left(\left(\frac{\varphi}{2}\right)^m\right),
\]
而定理 \ref{thm:xi-time-part70a-uniform-average-logdeg-two-state} 已证明
\[
\overline{\ell}_m
=
m\log\!\Bigl(\frac{2}{\varphi}\Bigr)
-\frac{F_m}{F_{m+2}}\log\varphi
+O\!\left(\left(\frac{\varphi}{2}\right)^m\right).
\]
相加便得到第一条公式。

再由 Binet 公式，
\[
\frac{|X_m|}{2^{m+1}}
=
\frac{F_{m+2}}{2^{m+1}}
=
\frac{\varphi^2}{2\sqrt5}\left(\frac{\varphi}{2}\right)^m
+O\!\bigl((2\varphi)^{-m}\bigr),
\]
且
\[
\frac{F_m}{F_{m+2}}=\varphi^{-2}+O(\varphi^{-2m}).
\]
用这两式乘第一条公式，得到
\[
g_m-\kappa_m+1
=
\frac{\varphi^2}{2\sqrt5}
\Bigl(
m\log\!\Bigl(\frac{2}{\varphi}\Bigr)
-\varphi^{-2}\log\varphi
+\log(2\pi)
\Bigr)
\left(\frac{\varphi}{2}\right)^m
+O\!\bigl(m(2\varphi)^{-m}\bigr)
+O\!\left(\left(\frac{\varphi}{2}\right)^{2m}\right).
\]
由于
\[
\frac{1}{2\varphi}<\left(\frac{\varphi}{2}\right)^2,
\]
故
\[
m(2\varphi)^{-m}=O\!\left(\left(\frac{\varphi}{2}\right)^{2m}\right).
\]
于是第二条公式成立。证毕。
\end{proof}

\noindent
由此，bin-fold 的两态一致渐近不但控制稳定层均匀经验谱与全部缩放矩，而且把均匀平均的对数退化度与规范体积修正的首项常数同时固定在同一末位二态账本之上。于是，均匀层二点律、幂和首项、平均对数势与规范体积差额最终被收束为同一条可显式求值的黄金比例主线。

\endinput
